% AVOI -- integrated manuscript-ready results text across all three stages.
% Stage 1 = original pre-registered experiment (historical, unchanged).
% Stage 2 = evidence-model adequacy safety condition + confirmation on new data (defined after Stage 1; not part of
%           the Stage-1 pre-registration).
% Stage 3 = adaptive next-evidence selection (WP1) and external-domain evidence-model validation (WP2), defined and
%           frozen after Stage 2 concluded. Stage 3 was NOT part of the original Stage-1 pre-registration and was NOT
%           anticipated by the Stage-2 protocol; it is reported here as a third, separately frozen stage of the same
%           research program, not retroactively folded into the original hypotheses.
\section{Results}

\subsection{Experimental history and design}
The evaluation proceeded in three stages, each frozen before its own data collection and each building only on data
and code artefacts explicitly designated development inputs for the next stage. Stage~1 was the original,
pre-registered experiment (hypotheses H1--H5), with no evidence-model adequacy condition. Its poor cross-testbed
transfer on AIT-LDS motivated an evidence-model adequacy safety condition, defined for Stage~2 -- this condition was
not part of the Stage-1 pre-registration. Stage~2 froze the adequacy policy, evaluated it on new IAM scenarios and
previously unused AIT-LDS/CERT partitions, and found only IAM evidence models eligible for autonomous use. Two
questions were left open by Stage~2's frozen design and could not have been asked within it: whether AVOI's
next-evidence choice is actually history-dependent rather than an artefact of a fixed initial ranking, and whether the
adequacy policy that blocked every non-IAM evidence model in Stage~2 was structurally unpassable or simply had not yet
been given a fair chance with correctly developed models. \textbf{Stage~3 was designed and frozen specifically to
answer these two questions after Stage~2's results were known; it is a new, third stage of this research program, not
a revision of the original pre-registration or a hypothesis the first two stages were testing for.} It reuses the
AVOI decision rule, loss matrix and adequacy policy unchanged, and treats every Stage-1/2 dataset as Stage-3
development data only -- never as new confirmation data.

\subsection{Stage 1}
AVOI substantially reduced the number of evidence acquisitions relative to exhaustive and fixed-order strategies on
all three original datasets. It did not uniformly reduce unsafe assurance relative to information gain (H2 not
supported), and exhaustive evidence collection achieved lower expected loss on the controlled IAM testbed. On
AIT-LDS, all methods showed high unsafe assurance ($\mathrm{UAR}\approx 0.31$), reflecting a cross-testbed evidence-model
transfer problem later addressed structurally in Stage~2.

\subsection{Stage 2: evidence-model adequacy}
Four of nine Stage-2 evidence models satisfied the predefined autonomous-assurance eligibility policy, all four on the
controlled IAM testbed; the AIT exfiltration model was classified UNSUPPORTED (inverted discrimination, AUROC 0.024),
the AIT intrusion and CERT models lacked sample support or were over-confident, and two further IAM models were
restricted by the calibration-slope rule. On the four validated IAM models, AVOI required 1.45 acquisitions per
confirmation episode versus 4.55 more for exhaustive collection and 3.16 more for a fixed checklist (Holm-adjusted
$p<10^{-128}$), with unsafe assurance identical to information gain. The gate produced no autonomous decisions for any
non-validated model and never converted model inadequacy into a false assurance, at the cost of coverage wherever it
acted (UAR 0.013$\to$0.002 across IAM, coverage 0.785$\to$0.451).

\subsection{Stage 3: is AVOI's evidence selection actually adaptive?}
Using IAM scenarios extended with a latent, AVOI-invisible branching mechanism (the same admissible evidence-action
set can require different next evidence depending on observed history, but the acquisition policy never sees the
latent mechanism directly), a group-aware permutation test rejected history-independence of AVOI's next-evidence
choice for all three tested objectives (Holm-adjusted $p=0.0003$ each), with 55--59\% of observed steps showing a
reversal relative to the policy's own initial ranking. \textbf{This state-dependent re-ranking is real}: it is not an
artefact of always following a fixed initial ranking. Whether this adaptivity reduces expected decision loss compared
with a policy that fixes its ranking after the first observation (``Frozen-Order AVOI'') was tested at matched
evidence budgets and was \textbf{not supported at the pre-registered, pooled level} (neither $k=2$ nor $k=3$ showed a
significant, correctly signed effect once pooled across the three tested objectives). Descriptively, adaptive
re-ranking helped significantly on one objective (loss lower by 0.086--0.113, Holm $p=0.0006$), made no measurable
difference on a second (the same evidence is collected regardless of order), and was numerically -- marginally
significantly at one horizon -- worse on the third, which was also the one objective whose evidence model failed the
adequacy calibration condition. We therefore do not claim that AVOI's re-ranking reduces expected decision loss in
general; we claim only that the re-ranking is real and that its value is objective- and model-adequacy-dependent.

\subsection{Stage 3: does the adequacy policy ever pass an external domain?}
Stage~2's AIT exfiltration failure was diagnosed as a non-monotone calibration-map artefact of a flawed two-testbed
development split (the raw discriminating statistic separates violations with AUROC 1.00 on every development
testbed), not a genuine, irreparable transfer failure. Using a corrected development procedure (calibrators restricted
to monotone maps; model selection by leave-one-development-testbed-out worst-fold AUROC) on three previously unused
AIT-LDS v2.0 testbeds and two previously unused CERT releases, one of three external claims -- CERT insider-threat
detection -- passed the unchanged Stage-2 adequacy policy on independent validation data (AUROC 0.895, Brier skill
score 0.744, expected calibration error 0.017; VALIDATED under all primary adequacy conditions). This is the first
non-IAM evidence model to pass the adequacy policy across both Stage~2 and Stage~3. The two AIT claims remained
RESTRICTED for insufficient sample support on the pre-committed target testbed; one of them (intrusion detection)
showed discrimination as strong and consistent as the validated CERT model (worst-development-fold AUROC 0.837,
confirmation-testbed AUROC 0.969) and was restricted on sample count alone, while the other (exfiltration detection)
showed genuinely weaker cross-testbed transfer (confirmation-testbed AUROC 0.729, falling to 0.563 on one testbed)
independent of its sample-support shortfall. On the validated CERT model, AVOI retained a significant acquisition
advantage over a fixed checklist at comparable autonomous coverage (2.42 vs.\ 4.85 mean acquisitions, Holm
$p<10^{-75}$), but this advantage is not claimed against exhaustive collection, whose higher coverage exceeded a
pre-registered no-trivial-abstention margin. The validated model itself had zero sensitivity to the violated class at
its tuned operating threshold, despite passing every adequacy condition -- a reminder that the adequacy policy
certifies probabilistic quality, not operational recall of the outcome that matters.

\subsection{What Stage 3 adds and does not add}
Stage 3 establishes that AVOI's adaptivity is genuine and that the adequacy policy is not structurally unpassable
outside the controlled IAM testbed -- both were open questions after Stage 2. It does not establish that adaptive
re-ranking is generally better than a fixed order, that AVOI generalizes to AIT-LDS, or that CERT-style evidence models
are reliable beyond the single validated claim and reduced release pool examined here. These qualifications are
reported alongside the positive results, following the same evidentiary standard applied throughout Stages 1 and 2.
